\origpage{34}{ROBERTSON AND SEYMOUR}
\noindent and $v_1, v_2 \in V(C)$, and there is a region $r \ne \Lambda$ of $H'$ incident with $v_1$ and $v_2$, and $F_1 \subseteq \Lambda$ and $F_2 \subseteq r$, where $F_1, F_2$ are the two components of $F - \{v_1, v_2\}$. From the hypothesis, $\bar{r} \cap \bar{\Lambda}$ is connected, and so there is a line $L \subseteq \bar{r} \cap \bar{\Lambda}$ with ends $v_1, v_2$. Let $\Delta_1 \subseteq \bar{\Lambda}$ be the closed disc bounded by the O-arc $F_1 \cup L$; and let $\Delta_2 \subseteq \bar{r}$ be the closed disc bounded by $F_2 \cup L$. (These exist because $\bar{\Lambda}$ is a closed disc, and $\bar{r}$ is simply-connected by (6.2).) Then $\Delta = \Delta_1 \cup \Delta_2$ is a closed disc bounded by $F$. Moreover, $H' \cap \Delta_1 = H' \cap L$, and $H' \cap \Delta_2 = H' \cap L$ by (6.2). Hence $H \cap \Delta = H \cap L$, as required. $\square$

Two drawings are \textit{isomorphic} if they are isomorphic as graphs.

\stmt
(6.4) \textit{Let $\Delta$ be a closed disc, let $H_1, H_2$ be drawings in $\Delta$ with $V(H_1) \cap \text{bd}(\Delta) = V(H_2) \cap \text{bd}(\Delta)$, and let $\phi \colon H_1 \to H_2$ be an isomorphism such that $\phi(v) = v$ for all $v \in V(H_1) \cap \text{bd}(\Delta)$. If $H_1$ is rigid in $\Delta$ then there is a homeomorphism $\psi \colon \Delta \to \Delta$ extending $\phi$.}

\textit{Proof.} Since $H_1$ is rigid, no vertex of $H_1$ in $\Delta - \text{bd}(\Delta)$ has valency $\le 1$; and we may assume that no vertex of $H_1$ in $\Delta - \text{bd}(\Delta)$ has valency $\le 2$, by suppressing such vertices. We may assume that $E(H_1) \ne \emptyset$ and $H_1$ is not a path, and so $|V(H_1) \cap \text{bd}(\Delta)| \ge 3$ since $H_1$ is rigid. Let $u, v \in V(H_1) \cap \text{bd}(\Delta)$ be consecutive in $\text{bd}(\Delta)$, and let $r_1$ be the region of $H_1$ in $\Delta$ including the part of $\text{bd}(\Delta)$ between $u$ and $v$. Define $r_2$ similarly, from $H_2$. If $u, v$ are not adjacent in $H_1$ we may add a new edge to $H_1$ and to $H_2$ joining $u, v$ drawn within the regions $r_1, r_2$ respectively; and the first enlarged drawing is still rigid, and still isomorphic to the second, and if the result holds for this pair then it holds for $H_1, H_2$. We may assume therefore that there is a circuit $C_1$ of $H_1$ with vertices the vertices of $V(H_1) \cap \text{bd}(\Delta)$, in order, and similarly for $C_2$. Let $r_1$ be as before. Then $r_1$ is incident with the edge joining $u, v$ because $H_1$ is rigid; and so, if we extend $\Delta$ to a sphere $\Sigma$, it follows that $U(C_1)$ is the boundary of the region $s_1$ of the drawing $H_1$ in $\Sigma$ which includes $\Sigma - \Delta$. Add a new vertex $v$ drawn in $s_1$, adjacent to every vertex of $C_1$, and let the enlarged graph be $J_1$; and define $J_2$ similarly. Then $J_1$ is a simple, 3-connected drawing in a sphere and $J_2$ is an isomorphic drawing, and the result follows from Whitney's theorem [7]. $\square$

\stmt
(6.5) \textit{Let $H$ be a rigid drawing in a closed disc $\Delta$. Let $H$ be a subgraph of $G$, let $\mathscr{S}$ be a 3-segregation of $G$ with $V(H) \cap \text{bd}(\Delta) \subseteq V(\mathscr{S})$, and let $\alpha$ be an arrangement of $\mathscr{S}$ in $\Delta$ such that $\alpha(v) = v$ for all $v \in V(H) \cap \text{bd}(\Delta)$. Then there is an arrangement $\alpha'$ of $\mathscr{S}$ in $\Delta$ such that $\alpha'$ is compatible with $H$.}

\textit{Proof.} First we observe the following.

\hspace*{12pt}(1) \textit{For every circuit $C$ of $H$ there are three disjoint paths of $H$ from $V(C)$ to $V(H) \cap \text{bd}(\Delta)$.}
